The following simulations where used in Chapter~\ref{chapter:simulation_finesse}. Section~\ref{sec:mode_match_to_arm} shows how the input beam parameter was determined so that single bounce simulations and simulations with arm cavities had the beam parameter of the input laser remained the same. Section~\ref{sec:sim_interferometer_params} contains tables of the spacings and properties of the optics for the LLO interferometer. Without a \texttt{cav} command, \textsc{Finesse} will default to using a 2mm beam with its waist at $z = 0$ with respect to the node at the laser.
\section{Mode Matching the Laser to the Arm Cavity}
\label{sec:mode_match_to_arm}
To find suitable numbers to use for the \texttt{gauss} command in the simulations of single-bounce and locked-arm measurements (see Section~\ref{sec:finesse_code} and Figure~\ref{fig:diagramofsim2}), the following test was performed. The main simulation from Section~\ref{sec:finesse_code} was modified so that it did not have an X arm or the SRC. The following beam parameters were found with a \texttt{bp} detector placed at the lasers node:
\begin{itemize}
	\item $w_0 = 0.00864064287177005$ m
	\item $z = 1339.83530090614$ m
\end{itemize}
To determine if the sign of z is correct, consider what happens if the beam parameter detector is moved 1m towards the cavity to \verb|bs_a|:
\begin{itemize}
	\item $w_0 = 0.00864064287177005 $
	\item $z = 1338.83530090614$ m
\end{itemize}
$z$ increasing suggests moving from \verb|n1| to \verb|bs_a| causes one to be closer to the waist. Intuitively, we know this is the case since the geometry of the cavity means the waist is within the cavity, and moving towards the cavity means you are moving towards the cavity waist. In other words, we know the beam should be getting smaller as we move from the laser to the cavity, which matches the behaviour seen here.
\par
However, setting the beam parameter at the laser with the first set of numbers gives \textit{terrible} mode matching. The simulation below investigates this. A photodiode which is masked from the TM00 mode 
\begin{verbatim}
pd HOM bs_d 
mask HOM 0 0 0
\end{verbatim}
was placed at \verb|bs_d| , which has the beam reflected from the arm cavity in it. The value of $z$ for the \texttt{gauss}  was varied and the power incident on \texttt{HOM}  was measured. To check that enough HOM were being used, the simulation was repeated with several \texttt{maxtem}  commands. The results had a quadratic fitted to them, and the minimum was taken as the optimal value. This is shown in Figure~\ref{fig:gausscommandfitting}. The values: 
\begin{itemize}
	\item $w_0 = 0.00864064287177005$ m
	\item $z = -1339.83683935$ m
\end{itemize}
give more than 99.99999\% mode matching, which is plenty for for this work. This is close to the original values that we obtained, except for a minus sign. This is due to the way \textsc{Finesse} computes the beam parameter. 
\begin{figure}
	\centering
	\includegraphics[width=0.9\linewidth]{sim_chapter_gauss_command_fitting}
	\caption{Optimisation of the parameters used in the \texttt{gauss} command to mode match the laser to the arms in the simulations described in Chapter~\ref{chapter:simulation_finesse} and Section~\ref{sec:finesse_code}. The blue, green and orange lines show how the mode matching changes as the beam waist location is changed. As the mode mismatch became larger, more HOM were required; this is why the blue and green/orange lines diverge. A quadratic was fitted to the data highlighted in red, and the minimum value was used as the waist location parameter for the \texttt{gauss} command.}
	\label{fig:gausscommandfitting}
\end{figure}
\clearpage
\section{Interferometer Parameters}
\label{sec:sim_interferometer_params}
The simulated interferometer uses nominal values for the LIGO arms cavities and the LLO SRC. The simulated arm cavity geometry is based on the design values to avoid contrast defects. The following sections contain tables specifying the design and measured values for each component of the simulation. 
\subsection{Radii of Curvature of optical components}
The optics in the SRC were simulated with the values from Table~\ref{tab:measured_vals_RoCs}. The arm cavity optics have been simulated using the values from Table~\ref{tab:desgin_vals_RoCs}.
\begin{table}[ht]
	\centering
	\begin{tabular}[h]{ c | c | c | c }
		Optic Name & RoC (m) & Focal length (m) & Reference \\
		\hline
		SRM &  $-5.6938$ &  $-2.8469$ &~\cite{Muzammil2016} \\
		SR2 &  $-6.427$ &  $-3.2135$ &~\cite{Muzammil2016} \\
		SR3 & $36$ &  $18$ &~\cite{Muzammil2016} \\
		ITM &  $1934$ &  $967.0$ &~\cite{Muzammil2016} \\
		ETM &  $2245$ &  $1122.5$ &~\cite{Muzammil2016} \\
	\end{tabular}
	\caption{The aLIGO design values of the RoC of the optics inside the SRC and the arms.}
	\label{tab:desgin_vals_RoCs}
\end{table}
\begin{table}[ht]
	\centering
	\begin{tabular}[h]{ c | c | c | c }
		Optic Name & RoC (m) & Focal length (m) & Reference \\
		\hline
		SRM-02 &  $-5.673 \pm 0.002$ & $-2.8365 \pm 0.001 $ &~\cite{galaxy_optics} \\
		SR2-04 &  $-6.406 \pm 0.006 $ &  $-3.203 \pm 0.003 $ &~\cite{SR2_04_measurement},\cite{galaxy_optics} \\
		SR3-01 & $35.97 \pm 0.012 $ &  $17.985 \pm 0.006$ &~\cite{SR3_01_figure_measurement},~\cite{galaxy_optics} \\
		ITM-08 &  $1938.44$ &  $969.22$  &~\cite{galaxy_optics} \\
		ETM-15 &  $2245.4$ & $1122.7$ &~\cite{galaxy_optics} \\
	\end{tabular}
	\caption{The measured values of the RoC of the optics at LLO that make up the SRC and the arms.}
	\label{tab:measured_vals_RoCs}
\end{table}
\subsection{Spacings}
The nominal values for the spacing between the optics at LHO/LLO, with a diagram, can be found in~\cite{aligo_layout_h1}. This information is also in~\cite{Muzammil2016}. The L1 Zemax parameters are found in~\cite{Zeemax_L1}. Table~\ref{tab:zeemax_spacings} contains the distances derived from the Zemax coordinates. Most spacings do not require sub mm accuracy, however SR3 to SR2 does since the Rayleigh range of this beam is about 4\,cm. The design value for the spacings are shown in Table~\ref{tab:design_spacings}. The values used in an earlier \textsc{Finesse} model of LLO~\cite{l1_finesse_2014} are shown in Table~\ref{tab:finesse_spacings}. The thickness of each optic is found in Table~\ref{tab:optic_thickness}.
\begin{table}[ht]
	\centering
	\begin{tabular}{ c | c }
		Space Name & Length (m) \\
		\hline
		EMTY to ITMY & $3994.4850$ \\
		ITMY to CP & $0.02$ \\
		CP to BS HR &  $4.8471$ \\
		BS AR to SR3 & $19.3659$ \\
		SR3 to SR2 & $15.4435$ \\
		SR2 to SRM & $15.7562$ \\
	\end{tabular}
\caption{Measured spacings between the optics at LLO\cite{Zeemax_L1}.}
\label{tab:zeemax_spacings}
\end{table}
\begin{table}[ht]
	\centering
	\begin{tabular}{ c | c }
		Space Name & Length (m) \\
		\hline
		EMTY to ITMY & $3994.5$ \\
		ITMY to CP & $0.05$ \\
		CP to BS HR & $4.8497$ \\
		BS AR to SR3 & $19.368$ \\
		SR3 to SR2 & $15.4607$ \\
		SR2 to SRM & $15.726$ \\
	\end{tabular}
	\caption{Design value for spacings\cite{aligo_layout_h1}.}
	\label{tab:design_spacings}
\end{table}
\begin{table}[ht]
	\centering
	\begin{tabular}{ c | c }
		Space Name & Length (m) \\
		\hline
		EMTY to ITMY & $3994.515$ \\
		ITMY to CP & $0.02$ \\
		CP to BS HR & $4.847$ \\
		BS AR to SR3 & $19.3661$ \\
		SR3 to SR2 & $15.4435$ \\
		SR2 to SRM &  $15.7566$ \\
	\end{tabular}
	\caption{Values used in an earlier \textsc{Finesse} model of LLO ~\cite{l1_finesse_2014}.}
	\label{tab:finesse_spacings}
\end{table}
\begin{table}[ht]
	\centering
	\begin{tabular}{ c | c }
		Optic & Thickness (m) \\
		\hline
		ITM & $0.200$ \\
		CP & $0.100$ \\
		BS & $0.0685$ \\
		Angled BS & $0.1202$ \\
	\end{tabular}
	\caption{Thickness of the optics between the ITM and the SRM.}
	\label{tab:optic_thickness}
\end{table}
\clearpage
\subsection{Parameters used in Main Simulation}
The parameters used in the main simulation were a combination of design values and measured values. The arm cavities were simulated using the design values to remove the contrast defect that would occur otherwise, and the SRC was simulated using measurements taken from LLO. The ITMs that were simulated have a thickness of 0.1\,m and a refractive index of 1.45. The thermal lens is implemented as shown in Figure~\ref{fig:diagramofsim2}. The parameters used are shown in Table~\ref{tab:nominal simulation}. 
\begin{table}[ht]
	\centering
	\begin{tabular}[h]{ c | c | p{45mm}  }
		Parameter & Value & Description \\
		\hline
		SRM RoC &  $-5.673\rm{m}$ &  Measured RoC of SRM-02 installed at LLO   \\
		SR2 focal length &  $-3.203\,\rm{m}$ &  Measured focal length of SR2-04 installed at LLO \\
		SR3 focal length & $17.985\,\rm{m}$ &  Measured focal length of SR3-01 installed at LLO \\
		ITM RoC &  $1934\,\rm{m}$ &  Design value \\
		ETM RoC &  $2245\,\rm{m}$ &  Design value \\
		ITM thermal lens &  $50\,\rm{km}$ & Design value \\
		ITM to ETM distance & $3994.5\,\rm{m}$& Design value \\
		ITM to BS Y &  $4.8471\,\rm{m}$ & As-built value for Y arm \\
		ITM to BS X &  $4.8471\,\rm{m}$ & As-built value for Y arm \\
		BS to SR3 & $19.3659\,\rm{m}$ & As-built value \\
		SR3 to SR2 & $15.4435\,\rm{m}$ & As-built value \\
		SR2 to SRM & $15.7562\,\rm{m}$ & As-built value \\
		ITM T & $0.0148$ & Power transmission\newline coefficient \\
		SRM T & $0.324$ & Power transmission \newline coefficient \\
		ITM L & $10\,\rm{ppm}$ & Scatter loss of ITM HR surface \\
		ETM L & $10\,\rm{ppm}$ & Scatter loss of ETM HR surface
	\end{tabular}
	\caption{Values for the major components of the simulation.}
	\label{tab:nominal simulation}
\end{table}
\clearpage
\section{\textsc{Finesse} Code}
This section contains the \textsc{Finesse} code used in Chapter~\ref{chapter:simulation_finesse}.
\label{sec:finesse_code}
\subsection{Michelson Interferometers with Thin Beamsplitter}
This \textsc{Finesse} code describes a Michelson interferometer where the beam splitter is implemented with a single \texttt{bs} command. This simulation will produce an incorrect result as the phase at the beam splitter is not properly accounted for.
\begin{lstlisting}[breaklines]
l l1 1 0 n1 
s sin 1 n1 na 

bs thin_bs 0.5 0.5 0 0 na nb nc nd 

s Yarm 1 nb nY
m Ymirror 1 0 0 nY dump

s Xarm 1 nc nX
m Xmirror 1 0 0 nX dump

pd AS nd

xaxis Xmirror phi lin 0 180 100
\end{lstlisting}
\subsection{Michelson Interferometers with Thick Beamsplitter}
This \textsc{Finesse} code describes a Michelson interferometer where the beam splitter is implemented with three \texttt{bs} and two \texttt{space} commands. This is the correct way to implement a beam splitter in \textsc{Finesse}.
\begin{lstlisting}[breaklines]
l l1 1 0 n1 
s sin 1 n1 na 

bs thick_bs_1 0.5 0.5 0 0 na nb nsub1 nsub3
s bs_subX 1p nsub1 nsub2
s bs_subY 1p nsub3 nsub4
bs thick_bs_X 0 1 0 0 nsub2 dump nc dump
bs thick_bs_AS 0 1 0 0 dump nsub4 dumo nd

s Yarm 1 nb nY
m Ymirror 1 0 0 nY dump

s Xarm 1 nc nX
m Xmirror 1 0 0 nX dump

pd AS nd

xaxis Xmirror phi lin 0 180 100
\end{lstlisting}
\subsection{Main Simulation}
The following \textsc{Finesse} code was used to simulate mode mismatches in the SRC.
\begin{lstlisting}[breaklines]
l l1 5M 0 n1
gauss beamParaml1 l1 n1 0.00864064287177005 -1339.83683935 

maxtem 12

s space1 1 n1 bs_a

#################################################
#bs beam_splitter 0.5 0.5 $phi_BS 0 bs_a bs_b bs_c bs_d

bs BSfront 0.5 0.5 0 0 bs_a bs_b nBSi1 nBSi3
s BSsubstrate1 1p 1.45 nBSi1 nBSi2
s BSsubstrate2 1p 1.45 nBSi3 nBSi4
bs BSback1 0 1 0 0 nBSi2 dump bs_c dump
bs BSback2 0 1 0 0 nBSi4 dump bs_d dump

##################################################
## Y ARM

s s_bs_to_ITMy 4.8471 bs_b nITMyARa

#ITMy
m1 ITMyAR 1 0 $phi_ITMy nITMyARa nITMyARb
s sITMy_thickness 0.05 1.45 nITMyARb nITMyTLa
lens ITMyTL 50k nITMyTLa nITMyTLb
s sITMy_thickness2 0.05 1.45 nITMyTLb nITMyHRa
m1 ITMy 0.0148 10u $phi_ITMy nITMyHRa nITMyHRb
attr ITMy Rc -1934

s space_arm_y 3994.5 nITMyHRb nETMy 

m1 ETMy 0 10u $phi_ETMy nETMy dump
attr ETMy Rc 2245

cav army ETMy nETMy ITMy nITMyHRb

##################################################

## X ARM

s s_bs_to_ITMx 4.8471 bs_c nITMxARa

#ITMx
m1 ITMxAR 1 0 $phi_ITMx nITMxARa nITMxARb
s sITMx_thickness 0.05 1.45 nITMxARb nITMxTLa
lens ITMxTL 50k nITMxTLa nITMxTLb
s sITMx_thickness2 0.05 1.45 nITMxTLb nITMxHRa
m1 ITMx 0.0148 10u $phi_ITMx nITMxHRa nITMxHRb
attr ITMx Rc -1934

s space_arm_x 3994.5 nITMxHRb nETMx

m1 ETMx 0 10u $phi_ETMx nETMx dump
attr ETMx Rc 2245

cav armx ETMx nETMx ITMx nITMxHRb

##################################################

## SRC 

s bs_to_SR3 19.3659 bs_d nSR3a

lens SR3 17.985 nSR3a nSR3b  

s SR3_to_SR2 15.4435 nSR3b nSR2a

lens SR2 -3.203 nSR2a nSR2b

s SR2_to_SRM 15.7562 nSR2b nSRMa

m1 SRM 0.324 0 0 nSRMa nSRMb
s SRM_substrate 0.07 1.45 nSRMb nSRMc
m SRM_AR 0 1 0 nSRMc nSRMd
attr SRM Rc -5.673

#cav src ITMx nITMxHRa SRM nSRMa

##################################################
#N.B these tunings don't represent how an IFO is kept on the dark fringe
const phi_ITMx 0 
const phi_ITMy 0
const phi_ETMx 0 
const phi_ETMy 0
const phi_BS 0  
##################################################

beam SRM_trans nSRMd
#mask SRM_trans 0 0 0

bp SRM_trans_w0 x w0 nSRMd
bp SRM_trans_w x w nSRMd
bp SRM_trans_r x r nSRMd

bp laser_w x w0 n1
bp laser_z x z n1

beam ITMy_refl nITMyHRa
beam ITMy_refl nITMyHRa

bp ITMy_refl_w x w nITMyHRa
bp ITMx_refl_w x w nITMxHRa

pd p_in bs_a*
pd p_ref bs_a

pd p_out nSRMd

pd p_out_hOM nSRMd
mask p_out_hOM 0 0 0

ad ad_HOM_2_y 2 0 0 nSRMd
ad ad_HOM_2_x 0 2 0 nSRMd

gouy src_gouy y sITMx_thickness sITMx_thickness2 s_bs_to_ITMx bs_to_SR3 SR3_to_SR2 SR2_to_SRM

trace 64
phase 2

xaxis SRM_trans x lin -10 10 200 
\end{lstlisting}

